% === Begin Label Data ===
% ID: 809
% 难度星级: 
% 标签: 
% 备注: 
% 组卷引用次数: 0
% === End  Label Data ===

\begin{problem}{2022}{G}{全国乙卷}{10}{概率}
某棋手与甲、乙、丙三位棋手各比赛一盘，各盘比赛结果相互独立．已知该棋手与甲、乙、丙比赛获胜的概率分别为 $ p_1, p_2, p_3 $，且 $ p_3 > p_2 > p_1 > 0 $．记该棋手连赢两盘的概率为 $ p $，则 (\hspace{1cm})
\begin{choices}
\choice{{ $ p $ 与该棋手和甲、乙、丙的比赛次序无关}}
\choice{{该棋手在第二盘与甲比赛时，$ p $ 最大}}
\choice{{该棋手在第二盘与乙比赛时，$ p $ 最大}}
\choice{{该棋手在第二盘与丙比赛时，$ p $ 最大}}
\end{choices}
\end{problem}

\begin{answer}
D
\end{answer}

\begin{solutions}
该棋手与甲、乙、丙三位棋手的比赛顺序共有 6 种可能：甲乙丙、甲丙乙、乙甲丙、乙丙甲、丙甲乙、丙乙甲，该棋手连赢两盘．可以是第一、第二盘连胜（此时第 3 盘胜负均可），或第一盘输但第二、第三盘连胜．

思路 1：直接计算．该棋手在第二盘与甲比赛时，考虑“乙甲丙”与考虑“丙甲乙”是等价的．此时，
$$
p = p_2 p_1 + (1 - p_2) p_1 p_3 \\
= p_1 (p_2 + p_3) - p_1 p_2 p_3,
$$
同理，该棋手在第二盘与乙比赛时，
$$
p = p_2 (p_1 + p_3) - p_1 p_2 p_3.
$$
该棋手在第二盘与丙比赛时，
$$
p = p_3 (p_1 + p_2) - p_1 p_2 p_3.
$$
显然，由 $ p_3 > p_2 > p_1 > 0 $ 可知，
$$
p_1 (p_2 + p_3) < p_2 (p_1 + p_3) < p_3 (p_1 + p_2).
$$
从而该棋手在第二盘与丙比赛时，$ p $ 最大，故正确选项为 D．

思路 2：特殊值法．为计算方便，可以把 $ p_1 $，$ p_2 $，$ p_3 $ 设成具体的数值．  
设 $ p_1 = 0.3 $，$ p_2 = 0.4 $，$ p_3 = 0.5 $，则该棋手在第二盘分别与甲、乙、丙比赛时，$ p $ 的值分别为 $ 0.42 $，$ 0.52 $，$ 0.58 $，可直接排除选项 B 和 C，且 $ p $ 的值与该棋手甲、乙、丙的比赛次序有关，进而排除选项 A，从而正确选项为 D．
\end{solutions}